\slajdid{07-s015}
\begin{frame}[label=07-s015]
\gusttekst\setlength{\parskip}{4pt}
\textbf{4. Netlist Language.} VHDL je moćan jezik sa kojim se unose novi dizajni na visokom nivou, ali je takođe koristan kao oblik komunikacije niskog nivoa između različitih alata u računarskom okruženju dizajna. Karakteristike strukturalnog jezika VHDL-a omogućavaju da se efikasno koristi kao jezik liste povezanosti mreža, zamenjujući (ili povećavajući) druge jezike liste povezanosti mreža kao što je EDIF. \par\vspace{4mm}
\textbf{5. Standardni jezik.} Jedan od najubedljivijih razloga da steknete iskustvo i znanje o VHDL-u je njegovo usvajanje kao standarda u zajednici elektronskog dizajna. Korišćenje standardnog jezika kao što je VHDL će praktično garantovati da nećete morati da bacate i ponovo preuzimate koncepte dizajna samo zato što metod unosa dizajna koji ste izabrali nije podržan u novijoj generaciji alata za dizajn. Korišćenje standardnog jezika takođe znači da je veća verovatnoća da ćete moći da iskoristite prednosti najsavremenijih alata za dizajn i da ćete imati pristup bazi znanja hiljada drugih inženjera, od kojih mnogi rešavaju slične probleme. \par\vspace{4mm}
\textbf{VHDL:} Jednostavno ili složeno koliko želite. Iako je istina da je VHDL veliki i složen jezik, zapravo nije teško početi sa njim. Koristite ga kako vam je potrebno i istražujte napredne funkcije dok postajete sigurniji.
\end{frame}
